\scanpage{21}\mspage{19}
On voit donc que $E$ doit satisfaire encore à une
condition supplémentaire, pour que \inserted{le} $\overline{\unsure{\mathscr{L}|\mathcal{L}}}_k$ correspondant
soit \unsure{isom.|iso.} à $\mathscr{L}_l$ i.e. pour que $E$ provienne d'un
foncteur $\varphi$. Je résume toutes les conditions

\[
\operatorname{Card}E=2(k-l)\tag*{a)}
\]
b) Les comp. conn. de $E$ qui ne contiennent pas
de flèche initiale ni finale sont de card. pair.

c) Si $\mathscr{L}_k$ et $\mathscr{L}_l$ ont $=$ direction initiale
($k$ et $l$ de $=$ parité), alors l'éventuelle
composante de $E$ contenant la flèche
initiale doit être elle aussi de cardinal pair
(\unsure{absence signifiant card de 0|absence signifiant card de .}), équivaut à la $=$
condition pour l'éventuelle comp. conn. contenant la flèche finale de $\mathscr{L}_k$). Si $\mathscr{L}_k$ et $\mathscr{L}_l$
ont des directions initiales différentes, il
faut que la flèche initiale de $\mathscr{L}_k$ appartienne
à une comp. conn. de cardinal \uline{impair} de
$E$ (ce qui équivaut à idem pour la flèche
finale).

D'autre part, le foncteur $\varphi$ est uniquement déterminé par
la connaissance de $E$, car \deleted{ils sont} il y a \unsure{une|un}
\uline{unique} iso. $\overline{\mathscr{L}}_k\xrightarrow{\sim}\mathscr{L}_l$ compatible aux sources et
buts, à cause du

\uline{Lemme} (immédiat) : \deleted{\unsure{tout|tel}} autom. de $\mathscr{L}_l$ qui
est l'identité sur l'\deleted{\unsure{point}} \inserted{objet} source (ou sur l'objet
but) est l'identité partout.

Les foncteurs $\varphi,\psi$ de tantôt correspondent
au choix, l'un de $E_\varphi=\{1,2,k-1,k\}$, l'autre de
$E_\psi=\{1,2,i,i+1\}$ \deleted{\unsure{tous les deux}} (avec $2$ composantes connexes
% archive sheet number: 20
